\documentclass[10pt,a4paper]{amsart}
\usepackage[margin=19mm]{geometry}
\usepackage{amsmath,amssymb,mathtools}
\usepackage[scr=rsfs,cal=euler]{mathalfa}
\usepackage{xfrac}
\usepackage[unicode,hidelinks,bookmarks=false]{hyperref}
\newcommand{\ie}{i.e.\ }
\newcommand{\cf}{cf.\ }
\newcommand{\resp}{respectively\ }
\newcommand{\Subsection}{\subsection}
\providecommand{\nobreakdash}{\nobreak\hbox{-}}
\setlength{\emergencystretch}{2em}
\multlinegap0pt
\usepackage{bm}
\usepackage{bbm}
\usepackage{dsfont}
\usepackage{xspace}
\usepackage{cancel}
\usepackage{mathtools}
\usepackage{amsthm}
\makeatletter
\@ifundefined{lemma}{\newtheorem{lemma}{Lemma}}{}
\makeatother
\title[Asymptotic Robustness in Biochemical Systems]{Asymptotic Robustness in Biochemical Systems}
\author{Hyukpyo Hong, Diego Rojas La Luz, Gheorghe Craciun}
\date{Source: arXiv:2505.21900v3 (CC BY 4.0)}
\begin{document}
\maketitle

\noindent\emph{Source and licence.} Asymptotic Robustness in Biochemical Systems. Version arXiv:2505.21900v3 (CC BY 4.0), distributed under the Creative Commons Attribution 4.0 International licence, as recorded in the arXiv licence metadata for this exact version and shown on the abstract page https://arxiv.org/abs/2505.21900v3. Full rights notice: licenses/arxiv-2505.21900-CC-BY-4.0.txt.\par\smallskip

\noindent\textit{Editorial scope.} Counted unit: the Lemma whose source label is \texttt{lem:LimitExists} in the article (source0.tex lines 566--574, source label \texttt{lem:LimitExists}), delivered together with the article's own complete original proof of it (source0.tex lines 576--608, one \texttt{proof} environment of 33 lines). No mathematics is written, repaired or re-derived: statement and proof are the article's own text line for line. Required context retained uncounted: none. Every definition, display and label of those lines is retained unchanged. Omitted from this package and not redistributed: 1--565, 575--575, 609--1050. Nothing inside a retained excerpt is omitted or altered: every retained body is delivered exactly as the article's lines are written, so no omission receipt of any kind accompanies this package. Citations: the retained excerpts carry no citation command at all, so no bibliography entry of the article is transcribed and no reference is redistributed. Reference adaptation: none was needed and none was made; every reference inside the delivered unit resolves inside it, so no unresolved reference or citation is produced. Printed slips kept literal: no printed word, symbol, hypothesis or number of the counted statement or of its proof was altered. Layout and class: the article's own class is \texttt{siamart250211} with packages including colortbl, mathtools, tikz, enumitem, chemfig, lipsum, amsfonts, amssymb, graphicx, epstopdf, algorithmic, amsopn; the wrapper is the lane's pinned \texttt{amsart} editorial wrapper at 10pt on A4, which restates the theorem-like environments the retained text needs and adds the article's own macro definitions that the retained text needs (none), with extra wrapper packages (bm, bbm, dsfont, xspace, cancel, mathtools). The delivered numbering is the unit's own. Assets: no retained span contains a figure environment, a graphics-inclusion command, a PDF-page inclusion, a table or a TikZ picture; no image file is copied into this package, the package declares no asset dependency and no image file is redistributed with it. Licence: version arXiv:2505.21900v3 of this article is distributed under the Creative Commons Attribution 4.0 International licence, as recorded in the arXiv licence metadata for this exact version and shown on the abstract page https://arxiv.org/abs/2505.21900v3. A full rights notice is delivered at licenses/arxiv-2505.21900-CC-BY-4.0.txt. Characters: every retained excerpt is pure ASCII, so the delivered engine, XeLaTeX with its default Latin Modern text font, reports no missing character.
\begin{lemma}\label{lem:LimitExists}
    Let $P(x,\lambda)$ be a polynomial with real variables $x$ and $\lambda$, and consider a continuous function $x(\lambda)$  defined for all sufficiently large $\lambda > \lambda_0$, satisfying $P(x(\lambda),\lambda)=0$. Then $\displaystyle \lim_{\lambda \to \infty} x(\lambda)$ exists (within the closed interval $[-\infty, \infty]$).

    Moreover, let us  write 
    \begin{equation}
        P(x,\lambda) = \lambda^m q(x) + r(x,\lambda),
    \end{equation}
    where $m$ is the highest degree of $\lambda$ in $P(x,\lambda)$, $q(x)$ is non-zero, and the degree of $\lambda$ in $r(x,\lambda)$ is less than $m$. Then, if the limit is finite, it is a root of $q(x)$.
\end{lemma}

\begin{proof}[Proof of Lemma~\ref{lem:LimitExists}]
    To show that the limit exists it is enough to prove that
    \begin{equation}
        \liminf_{\lambda\to\infty} x(\lambda) = \limsup_{\lambda\to\infty} x(\lambda).
    \end{equation}
    Write $P(x,\lambda)$ as
    \begin{equation}
        P(x,\lambda) = \lambda^m q(x) + r(x,\lambda),
    \end{equation}
    where $m$ is the highest degree of $\lambda$ in $P(x,\lambda)$, $q(x)$ is non-zero, and the degree of $\lambda$ in $r(x,\lambda)$ is less than $m$. 
    
    If $m=0$, then $P(x,\lambda)=P(x)=q(x)$, which have only finitely many roots. Since $x(\lambda)$ is continuous and must admit a value which is a root of $P(x)=0$, it must be a constant function. Thus in this case the limit exists and it is a root of $q(x)$.
    % \color{black}
    
    We now assume that $m>0$. Then, 
    \begin{align}
        \lambda^m q(x(\lambda)) + r(x(\lambda),\lambda) &= 0,\\
        \Rightarrow q(x(\lambda)) + \dfrac{1}{\lambda^m} r(x(\lambda),\lambda) &= 0.
    \end{align}
    Suppose by contradiction that 
    $\liminf_{\lambda\to\infty}x(\lambda) = c_1$ is strictly less than \\
    $\limsup_{\lambda\to\infty}x(\lambda) = c_2$, with $c_1,c_2\in[-\infty,\infty]$. Then, for all $c \in (c_1,c_2)$, there exists a sequence $\lambda_k \to \infty$ such that
    \begin{equation}
        \lim_{k\to\infty} x(\lambda_k) = c
    \end{equation}
    by the Intermediate Value Theorem. 
    Then, in the limit $\dfrac{1}{\lambda^m} r(x(\lambda),\lambda) = 0$ because $x(\lambda)$ is bounded.
    But then, as $q(x)$ is a polynomial, it is continuous, so 
    \begin{equation}
        \lim_{k\to\infty} q(x(\lambda_k)) = q(c) = 0,
    \end{equation}
    which implies $c$ is a root of $q$ for all $c\in(c_1,c_2)$. Since the non-zero polynomial $q(x)$ can have only finitely many roots, it leads to a contradiction. Thus, the limit exists, and if it is finite, it must be a root of $q(x)$.
\end{proof}

\end{document}
